% !TEX TS-program = pdflatex
\documentclass[final]{beamer}
\usepackage[size=a0,orientation=landscape,scale=1.35]{beamerposter}

\usepackage[ngerman]{babel}
\usepackage[T1]{fontenc}
\usepackage[utf8]{inputenc}
\usepackage{amsmath,amssymb,booktabs}
\usepackage{graphicx}
\usepackage{tabularx}
\usepackage{ragged2e}
\usepackage{helvet}
\renewcommand{\familydefault}{\sfdefault}
\usepackage{xcolor}
\definecolor{hsblue}{RGB}{0,82,155}
\setbeamercolor{title}{fg=hsblue}
\setbeamercolor{block title}{fg=white,bg=hsblue}
\setbeamercolor{block body}{fg=black,bg=white}

\title{\bfseries Sind Rechnungen am Abend höher als mittags?}
\author{\textbf{Team Statistik 1} (B.Sc. Data Science, 1.\ Semester)}
\institute{\textbf{Technische Hochschule Ingolstadt – Fakultät Wirtschaftsingenieurwesen}}
\date{\today}

\newcommand{\smallgap}{\vspace{0.6cm}}

\begin{document}
\begin{frame}[t]

% Kopf
\begin{columns}[T]
  \begin{column}{0.98\textwidth}
    \vskip1cm
    \begin{center}
      {\usebeamercolor[fg]{title}\Huge\inserttitle}\\[0.6cm]
      {\Large\insertauthor}\\[0.25cm]
      {\large\insertinstitute}\\[0.6cm]
      \rule{\textwidth}{1pt}
    \end{center}
  \end{column}
\end{columns}

\vspace{1cm}

% Hauptinhalt
\begin{columns}[T,totalwidth=\textwidth]

% Spalte 1
\begin{column}{0.32\textwidth}
  \begin{block}{1. Forschungsfrage}
    In vielen Restaurants wird vermutet, dass Gäste am Abend höhere Rechnungen bezahlen als mittags.  
    Doch ist dieser Unterschied auch \emph{statistisch nachweisbar}?  
    Wir prüfen diese Frage anhand des bekannten \texttt{tips}-Datensatzes aus Python.  

    \smallgap
    \textbf{Forschungsfrage:}  
    Gibt es einen signifikanten Unterschied der durchschnittlichen \texttt{total\_bill} zwischen \textbf{Lunch} und \textbf{Dinner}?  

    \smallgap
    \textbf{Hypothesen:}
    \[
      H_0: \mu_{\text{Lunch}} = \mu_{\text{Dinner}} 
      \qquad \text{vs.} \qquad 
      H_1: \mu_{\text{Lunch}} \neq \mu_{\text{Dinner}}
    \]

    \smallgap
    Es handelt sich also um einen klassischen \textbf{Gruppenvergleich} mit einer \textbf{kategorialen} und einer \textbf{stetigen} Variable.
  \end{block}

  \begin{block}{2. Datenbasis und Variablen}
    Der Datensatz \texttt{tips} enthält Restaurantrechnungen, Trinkgelder und Zusatzinformationen.  
    Für diese Analyse betrachten wir nur:
    \begin{itemize}
      \item \texttt{time} — Kategorie: \textit{Lunch} oder \textit{Dinner}
      \item \texttt{total\_bill} — Höhe der Gesamtrechnung in US-Dollar
    \end{itemize}
    Insgesamt liegen \textbf{244 Beobachtungen} vor, davon 68 zu Lunch und 176 zu Dinner.  

    \smallgap
    \textbf{Ziel:} Überprüfen, ob sich die Mittelwerte signifikant unterscheiden.

    \smallgap
    \centering
    \includegraphics[width=0.95\linewidth]{countplot.pdf}\\[-0.1cm]
    \small Verteilung der Beobachtungen: mehr Rechnungen abends als mittags
    \normalsize
  \end{block}
\end{column}

% Spalte 2
\begin{column}{0.34\textwidth}
  \begin{block}{3. Deskriptive Statistik}
    \centering
    \includegraphics[width=0.95\linewidth]{histogram.pdf}\\[-0.1cm]
    \small \textbf{Histogramm:} Verteilung der Rechnungsbeträge für Lunch und Dinner
    \normalsize

    \smallgap
    \textbf{Zentrale Kennzahlen:}
    \begin{center}
    \begin{tabularx}{0.9\linewidth}{Xccc}
      \toprule
      & $\bar{x}$ & $s$ & $n$ \\
      \midrule
      Lunch  & 17.17 & 7.71 & 68 \\
      Dinner & 20.80 & 9.14 & 176 \\
      \bottomrule
    \end{tabularx}
    \end{center}

    \textit{Medians:} Lunch = 15.97, Dinner = 18.39 \\
    \textit{Min–Max:} Lunch = [7.51; 43.11], Dinner = [3.07; 50.81]

    \smallgap
    Bereits visuell zeigt sich: Dinner-Rechnungen sind tendenziell höher und breiter gestreut.
  \end{block}

  \begin{block}{4. Annahmencheck}
  Bevor der t-Test eingesetzt wird, müssen einige Voraussetzungen überprüft werden:

  \textbf{1. Normalverteilung:}  
  Der Shapiro-Wilk-Test ergibt:
  \begin{itemize}
    \item Lunch: $W = 0{,}869$, $p = 3{,}6\cdot10^{-6}$, Schiefe = 1.48, Exzess = 2.21
    \item Dinner: $W = 0{,}933$, $p = 2{,}7\cdot10^{-7}$, Schiefe = 1.03, Exzess = 1.04
  \end{itemize}
  Die Daten weichen also formal von der Normalverteilung ab.  
  Da beide Stichproben aber hinreichend groß sind ($n>30$), gilt der t-Test als robust gegen moderate Abweichungen.  

  \smallgap
  \textbf{2. Varianzgleichheit:}  
  Levene-Test (Median) ergibt $p = 0{,}102$ → kein signifikanter Unterschied der Varianzen.  
  Wir können daher den klassischen t-Test verwenden.
  \end{block}

\end{column}

% Spalte 3
\begin{column}{0.32\textwidth}
  \begin{block}{5. Testdurchführung und Ergebnisse}
    Für den Vergleich der Mittelwerte wurde ein \textbf{t-Test für unabhängige Stichproben} durchgeführt:

    \[
      t = -2.90, \quad p = 0.0041
    \]

    Damit ist der Unterschied \textbf{statistisch signifikant} auf dem 5\%-Niveau.  
    Es ist also sehr unwahrscheinlich, dass der beobachtete Unterschied rein zufällig auftritt.

    \smallgap
    Zur Bewertung der praktischen Relevanz wurde zusätzlich die \textbf{Effektstärke (Cohen’s d)} berechnet:

    \[
      d = -0.41
    \]
    Nach der Konvention von Cohen entspricht dies einem \textbf{mittleren Effekt} — also einem spürbaren, aber nicht extremen Unterschied.

    \smallgap
    \centering
    \includegraphics[width=0.95\linewidth]{boxplot.pdf}\\[-0.1cm]
    \small Boxplot der Rechnungsbeträge für Lunch und Dinner
    \normalsize

    \smallgap
    Im Mittel liegen die Dinner-Rechnungen etwa \textbf{3,6 \$ über den Lunch-Rechnungen}.
  \end{block}

  \begin{block}{6. Fazit und Ausblick}
    \textbf{Fazit:}  
    Gäste geben am Abend tatsächlich im Durchschnitt mehr Geld aus als mittags.  
    Der Unterschied ist \textbf{statistisch signifikant} und zeigt einen \textbf{moderat ausgeprägten Effekt}.  
    Trotz formaler Abweichung von der Normalverteilung sind die Ergebnisse aufgrund der robusten Stichprobengröße belastbar.

    \smallgap
    \textbf{Limitationen:}
    \begin{itemize}
      \item Beobachtungsdaten – keine Kausalinterpretation möglich
      \item mögliche Störfaktoren (z.\,B. Gruppengröße, Alkohol, Wochentag)
    \end{itemize}

    \textbf{Ausblick:}
    \begin{itemize}
      \item Erweiterung: Zusammenhang zwischen Rechnung und Trinkgeld
      \item multivariate Modelle zur Kontrolle weiterer Variablen
    \end{itemize}

  \end{block}
\end{column}

\end{columns}

% Fuß
\begin{columns}[T]
  \begin{column}{0.98\textwidth}
    \begin{center}
      \rule{\textwidth}{1pt}\\[0.35cm]
      {\small Statistik I – Gruppenvergleich \quad|\quad Datensatz: \texttt{tips} \quad|\quad THI, WS 2025/26}
    \end{center}
  \end{column}
\end{columns}

\end{frame}
\end{document}
